% =====================================================================
%  CHAPTER 8: 儿童少年慢性病（对应大纲第八章）
%  内容来源：往届笔记第十章
% =====================================================================
\chapter{儿童少年慢性病}
\setcounter{chapter}{8}
\chapterintro{
\textbf{掌握：}儿童少年常见慢性病，包括哮喘、糖尿病、高血压和恶性肿瘤。
}

% ===== 掌握内容 =====
\section{儿童少年哮喘}

\begin{defbox}
\textbf{哮喘（asthma）：}一种以肥大细胞反应、嗜酸性粒细胞和T细胞浸润为主，伴淋巴细胞、巨噬细胞等参与调节的气道慢性炎症性疾病。临床表现为突然而反复发作的喘息、气短、胸闷和咳嗽，同时有可变性呼气气流受限，多于运动后或夜间和凌晨发作。
\end{defbox}

\subsection{分类}

\begin{keybox}
\textbf{按病因分类：}

\textbf{1. 外源性哮喘/特应性哮喘/季节性哮喘}
\begin{enumerate}[leftmargin=*]
    \item 发生于特应性人群，有其它过敏性疾病史。
    \item 暴露于较高水平的环境变应原。
    \item 发作有明显季节性，一般儿童青少年多见，常有哮喘家族史。
\end{enumerate}

\textbf{2. 内源性哮喘}
\begin{enumerate}[leftmargin=*]
    \item 发生于非特应性人群，一般无过敏性疾病史、明显季节性与哮喘家族史。
    \item 运动和病毒感染是最常见诱发因素。
    \item 发病年龄较晚，女性多见。
\end{enumerate}
\end{keybox}

\subsection{病因和影响因素}

\begin{keybox}
\textbf{1. 遗传因素}

\textbf{2. 环境因素}

（1）\textbf{变应原：}引发过敏反应的物质。
\begin{itemize}
    \item 室内变应原：尘土、尘螨、宠物唾液、皮毛、病毒、细菌、真菌。
    \item 室外变应原：花粉、真菌。
\end{itemize}

（2）\textbf{呼吸道病毒感染：}损伤气道黏膜，引发炎症反应。
\begin{itemize}
    \item 婴幼儿：呼吸道合胞病毒、副流感病毒、腺病毒。
    \item 学龄期：鼻病毒、流感病毒、副流感病毒、肺炎支原体。
\end{itemize}

（3）\textbf{气候变化：}气温、气湿、气压，雷暴、沙尘暴。

（4）\textbf{空气污染：}工业烟雾、光化学烟雾、装修材料（甲醛）、烹调油烟。

（5）\textbf{食物：}
\begin{itemize}
    \item 动物蛋白：鸡蛋、牛奶、鱼虾
    \item 坚果类：花生、腰果
    \item 谷类：小麦、玉米
    \item 食品添加剂：亚硝酸盐、防腐剂
\end{itemize}

（6）\textbf{其他因素：}首饰、香水；微波炉、电视、电脑、音响产生的电磁辐射；环境内分泌干扰物；肥胖症。
\end{keybox}

\subsection{预防控制}

\begin{keybox}
\textbf{1. 三级预防}
\begin{enumerate}[leftmargin=*]
    \item \textbf{一级预防：}避免接触可致敏性物质，消除一切可导致哮喘发生发展的高危致病因素；从母孕期入手，重视妊娠环境。
    \item \textbf{二级预防：}对出现变应性疾病初发症状患儿加强预防，防止再次发作；对患特异性皮炎、变应性鼻炎者加强观察，警惕出现哮喘并发症；对伴喘息等哮喘早期症状者，尽量减轻症状，减少发作次数。
    \item \textbf{三级预防：}防止症状呈不可逆性，避免出现后遗症。
\end{enumerate}

\textbf{2. 避免接触变应原}
\begin{enumerate}[leftmargin=*]
    \item 减少室内变应原
    \item 避免接触室外变应原
    \item 减少与呼吸道感染患者接触
\end{enumerate}

\textbf{3. 生活规律，适度锻炼}
\begin{enumerate}[leftmargin=*]
    \item 生活制度规律是对哮喘过敏儿最积极主动的预防措施。
    \item 避免剧烈活动和过度疲劳；适宜适度的体育锻炼可有效提高身体素质，降低特异性体质因素；剧烈运动时需进行预防性用药。
\end{enumerate}

\textbf{4. 普及哮喘防控知识}
\begin{enumerate}[leftmargin=*]
    \item 不仅应针对儿童和家长，还应针对周围相关人群。
    \item 通过宣传教育，提高家长和患儿对哮喘防治的信心，提高治疗依从性和主观能动性，控制各种触发因素，巩固治疗效果，提高生活质量。
\end{enumerate}
\end{keybox}

\section{儿童少年糖尿病}

\begin{defbox}
\textbf{糖尿病（diabetes mellitus, DM）：}一组由遗传、环境因素交互作用导致的慢性并以血糖升高为主要特征的临床综合征；因胰岛素分泌量绝对或相对不足及靶组织细胞对胰岛素敏感性降低，引发糖、蛋白质、脂肪、电解质和水等一系列代谢紊乱。
\end{defbox}

\subsection{分类}

\begin{keybox}
\textbf{1. 1型糖尿病（T1DM）/ 胰岛素依赖型糖尿病：}胰岛细胞受损导致体内胰岛素分泌绝对缺乏，从而引发系统性代谢紊乱。临床特点：起病急，多食、多尿、多饮、体重减轻。

\textbf{2. 2型糖尿病（T2DM）/ 非胰岛素依赖型糖尿病：}胰岛素抵抗为主，伴胰岛素分泌不足，导致机体调节葡萄糖代谢能力下降的代谢疾病。临床特点：起病缓慢、隐匿，临床症状相对较轻，有较强家族聚集性。
\end{keybox}

\subsection{病因和影响因素}

\begin{keybox}
\textbf{1. 遗传因素：}约25-50\%糖尿病患者有明显家族史，T1DM遗传倾向比T2DM更明显。

\textbf{2. 环境因素：}
\begin{enumerate}[leftmargin=*]
    \item 不良饮食行为：与大量摄入高热量、高糖、高脂、缺乏纤维素的饮食结构和方式直接相关。
    \item 肥胖：公认的最重要诱因之一。
    \item 体力活动不足。
    \item 心理社会因素。
\end{enumerate}
\end{keybox}

\subsection{预防控制}

\begin{keybox}
以往将儿童期糖尿病防治重点放在T1DM上，但近年来随着儿童少年中T2DM发病率上升，积极开展T2DM预防已成为紧迫而重要的任务。

\begin{enumerate}[leftmargin=*]
    \item \textbf{监测易感人群。}
    \item \textbf{开展生活方式干预，积极防治肥胖：}
    \begin{enumerate}[leftmargin=*]
        \item 合理膳食：低盐、低糖、低脂、高纤维、维生素充足。
        \item 培养科学生活方式：改变不良饮食习惯，不吸烟、不酗酒，积极参加体育锻炼。
        \item 定期检测体重：已发生超重者，从科学膳食、有氧锻炼、合理生活制度、建立健康行为等方面采取综合措施控制体重；已发展为肥胖者应在医师指导下，采取综合措施稳步降低体重。
    \end{enumerate}
    \item \textbf{培养运动习惯：}每天坚持≥1h中等强度运动，改善糖、脂肪代谢。
\end{enumerate}
\end{keybox}

\section{儿童少年高血压}

\subsection{定义与分类}

\begin{defbox}
\textbf{高血压（hypertension）：}以体循环动脉血压（收缩压和/或舒张压）增高为主要特征一类综合征。成人以收缩压≥140mmHg、舒张压≥90mmHg为高血压划界值，可伴心、脑、肾等器官功能或器质性损害的临床综合征。

\textbf{儿童高血压：}平均收缩压（SPB）和/或平均舒张压（DBP）等于或高于同年龄、同性别、同身高儿童的第95百分位数（P95）。
\end{defbox}

\begin{keybox}
\textbf{按病因分类：}

\textbf{1. 原发性高血压（primary hypertension）：}以血压升高为主要临床表现，伴或不伴多种心血管危险因素的综合征，约占患者总数95\%以上。在儿童晚期和青少年中较为常见，一般很少出现临床症状，多在常规检查时被发现。

\textbf{2. 继发性高血压（secondary hypertension）：}继发于其他疾病（如肾小球肾炎、主动脉缩窄等）。最常见的是由肾脏及肾上腺疾病所致或内分泌性高血压。高血压仅是原发性疾病临床表现之一，血压可暂时性或持久性升高。
\end{keybox}

\subsection{病因和影响因素}

\begin{keybox}
\textbf{1. 遗传因素：}明显家族聚集性。

\textbf{2. 环境因素：}
\begin{enumerate}[leftmargin=*]
    \item \textbf{膳食因素：}血压与能量、糖、蛋白质、脂肪、胆固醇和钠盐（相关性最突出）摄入呈正相关，与钾摄入呈负相关。
    \item \textbf{体力活动：}保护性因素。
    \item \textbf{超重肥胖：}重要危险因素。
\end{enumerate}

\textbf{3. 生长发育进程}
\begin{enumerate}[leftmargin=*]
    \item 从出生开始，血压随年龄增长而增高，成年后基本稳定。
    \item \textbf{轨迹现象：}对儿童血压进行长期跟踪发现，尽管血压随年龄变化，但多数儿童血压在同年龄组中所处百分位数大致保持不变。→儿童时期血压状况可能对成年后血压水平产生一定影响，也为早期关注和干预儿童血压提供了一定依据。
\end{enumerate}
\end{keybox}

\subsection{预防控制}

\begin{keybox}
\textbf{1. 群体预防——一般性预防}
\begin{enumerate}[leftmargin=*]
    \item 健康教育：普及高血压危害和预防知识。
    \item 学校-社区-家庭联动的预防高血压活动：倡导合理饮食、积极锻炼、合理生活作息制度和建立健康生活方式。
    \item 学校体育运动：控制体重，科学减肥。
    \item 建立健康行为：戒除吸烟、酗酒等不健康行为。
    \item 健全学生体检制度：每学期至少测1次血压。
\end{enumerate}

\textbf{2. 高危人群预防——特殊预防}
\begin{enumerate}[leftmargin=*]
    \item \textbf{原发性高血压儿童少年：}血压持续超过筛检标准P95并排除继发性原因者。原则上不考虑使用药物治疗，着重采取改变生活方式的治疗方法，定期复查血压。
    \item \textbf{高血压倾向儿童少年：}血压虽未超过筛检标准但处于同性别-年龄组血压高位（P80-P95）者。定期复查血压，开展行为干预。
    \item \textbf{高血压易患儿童少年：}有高血压家族史者。培养健康生活方式，定期复查血压。
\end{enumerate}
\end{keybox}

\section{儿童少年恶性肿瘤}

\begin{defbox}
\textbf{肿瘤（tumor）：}机体在各种致癌因素作用下，局部组织细胞在基因水平上失去对生长的正常调控，细胞异常增生形成的新生物，分为良性肿瘤和恶性肿瘤。
\end{defbox}

\begin{keybox}
\textbf{恶性肿瘤的危险因素：}

\begin{enumerate}[leftmargin=*]
    \item \textbf{环境因素：}（1）化学因素（最重要）：烷化剂、多环芳烃；（2）物理因素：辐射、紫外线；（3）生物因素：病毒、霉菌、寄生虫。
    \item \textbf{遗传因素。}
    \item \textbf{生活方式与行为因素：}吸烟、饮酒、饮食（膳食结构不合理）、不良生活习惯（长期熬夜）。
    \item \textbf{个体心理特征：}C型行为。
\end{enumerate}

\textbf{预防控制：}

\begin{enumerate}[leftmargin=*]
    \item \textbf{一级预防（病因预防）：}（1）保护、改善环境，防止和消除环境污染；（2）改变与肿瘤发生有关的行为危险因素；（3）少吃过烫、过辣食物，进食细嚼慢咽，防止食道因反复损伤诱发癌变；（4）儿童期开始加强防癌健康教育，提高自我保健能力。
    \item \textbf{二级预防：}早发现、早诊断、早治疗。
    \item \textbf{三级预防：}（1）减少肿瘤并发症；（2）防止伤残，提高生存率；（3）提供临终关怀，缓解晚期病人痛苦，提高生存质量。
\end{enumerate}
\end{keybox}

% ----- 复习题 -----
\reviewcontent{
\section{复习题}

\begin{enumerate}[leftmargin=*, itemsep=8pt]
    \item \textbf{【名词解释】}哮喘；糖尿病；高血压；儿童高血压；肿瘤；轨迹现象（血压）
    \item \textbf{【简答题】}简述哮喘的分类（外源性和内源性）及主要区别。
    \item \textbf{【简答题】}简述哮喘的环境影响因素和三级预防策略。
    \item \textbf{【简答题】}比较1型糖尿病和2型糖尿病的主要区别。
    \item \textbf{【简答题】}简述儿童高血压的分类及高危人群的三级预防策略。
    \item \textbf{【简答题】}简述恶性肿瘤的危险因素和三级预防策略。
\end{enumerate}

\subsection*{参考答案要点}

\begin{enumerate}[leftmargin=*, itemsep=5pt]
    \item \textbf{哮喘}：以肥大细胞反应、嗜酸性粒细胞和T细胞浸润为主的气道慢性炎症性疾病。\textbf{糖尿病}：由遗传、环境因素交互作用导致的以血糖升高为主要特征的临床综合征。\textbf{高血压}：以体循环动脉血压增高为主要特征一类综合征。\textbf{儿童高血压}：平均收缩压和/或舒张压等于或高于同年龄、同性别、同身高儿童P95。\textbf{轨迹现象（血压）}：尽管血压随年龄变化，但多数儿童血压在同年龄组中所处百分位数大致保持不变。

    \item 外源性哮喘：发生于特应性人群，有过敏性疾病史，有明显季节性，儿童青少年多见，常有家族史。内源性哮喘：非特应性人群，无过敏史、季节性与家族史，运动和病毒感染最常见诱因，发病年龄较晚，女性多见。

    \item 环境因素：变应原（室内：尘土/尘螨/宠物皮毛；室外：花粉/真菌）；呼吸道病毒感染；气候变化；空气污染；食物（动物蛋白/坚果/谷类/食品添加剂）；其他（电磁辐射/EDCs/肥胖症）。三级预防：一级（避免接触致敏原，从母孕期入手）；二级（防止再次发作，警惕并发症）；三级（防止症状不可逆，避免后遗症）。生活规律、适度锻炼、普及防控知识。

    \item T1DM：胰岛素依赖型，胰岛细胞受损导致胰岛素绝对缺乏，起病急，多食多尿多饮体重减轻。T2DM：非胰岛素依赖型，胰岛素抵抗为主伴分泌不足，起病缓慢隐匿，症状较轻，家族聚集性强。遗传倾向T1DM比T2DM更明显。

    \item 分类：原发性高血压（占95\%以上，儿童晚期和青少年常见）；继发性高血压（继发于其他疾病，最常见肾脏及肾上腺疾病）。三级预防：(1)原发性→不用药，改变生活方式，定期复查；(2)倾向者（P80-P95）→定期复查，行为干预；(3)易患者（有家族史）→培养健康生活方式，定期复查。

    \item 危险因素：(1)环境因素（化学因素最重要：烷化剂/多环芳烃；物理：辐射/紫外线；生物：病毒/霉菌/寄生虫）；(2)遗传因素；(3)生活方式与行为（吸烟/饮酒/饮食/熬夜）；(4)个体心理特征（C型行为）。预防：一级（病因预防，改善环境，改变行为危险因素，加强防癌健康教育）；二级（早发现早诊断早治疗）；三级（减少并发症，防止伤残，提供临终关怀）。
\end{enumerate}
}

\newpage
